% ---------------------------------------------------------------------------
\section{Data}
\label{sec:data}

\subsection{The Voynich manuscript}

We use three transliterations in IVTFF format \citep{zandbergen_ivtff,zandbergen_transliterations}: Zandbergen--Landini
(ZL3b, EVA alphabet; our reference), Takahashi (IT2a, basic EVA) and Glen Claston (GC2a, v101 alphabet, which segments
some glyphs differently). They are three largely independent readings, but they share the locus scheme curated by
Zandbergen: line segmentation and page metadata are the same in all three. Replicating a line-level result on the three
files therefore does not replicate where a line ends.

\paragraph{Treatment of the transliteration} We use the running paragraph text: 207 pages, about 4{,}100 lines and
34{,}900 words; labels, circular and radial text enter only the tests that say so. Uncertain readings
(\texttt{[a:b]}) take the first option; ligature braces are removed; rare glyphs (\texttt{@nnn;}) and illegible glyphs
(\texttt{?}) make a word unusable, and such words (0.6\%) are dropped; depending on the experiment, the pair across a
dropped word is either skipped or counted as adjacent. The uncertain space (EVA comma) counts as a space unless stated; the tests on
uncertain spaces treat it separately. The EVA sequences \eva{ch}, \eva{sh}, \eva{cth}, \eva{ckh}, \eva{cph} and
\eva{cfh} count as single glyphs. Pages, bifolia and quires are taken from the IVTFF page headers: 50 bifolia and 16
quires carry text. Hands follow \citet{davis2020} and languages follow \citet{currier1976}, both as recorded in the files
page by page.

\subsection{Comparison texts}

\begin{itemize}
\item \textbf{Natural and constructed languages.} The publicly released texts of \citet{gaskellbowern2022}: 71 texts,
  of which 56 are in 24 natural languages (usually a New Testament and a technical text per language, historical and
  modern; Latin also in an abbreviated form) and 15 in eight constructed languages (Esperanto, Interlingua, Klingon,
  Lojban, LOLCat, Neo-Quenya, Toki Pona, Volapük; LOLCat is a register of English rather than a constructed language).
  Eleven of the technical texts are translations of the Wikipedia article on the manuscript; they contain at most one
  or two EVA words each. The collection also provides the texts we cite by name (the Qur'an, Avicenna, the \emph{Hatton
  Gospels}, Alberti's \emph{Della Pittura}, Pliny). In these files the lines are those of the source editions (median eight
  words; about a quarter end with punctuation, and verses or paragraphs are kept apart), not the lines of manuscripts. Comparisons use the first 10{,}000 words of each text and, since
  version~2, one value per natural language (the median of its texts), with constructed languages reported apart
  (experiment \expcode{e3c84}).
\item \textbf{Other texts:} the Bible in about 100 languages \citep{christodouloupoulos2015}; Latin texts of the Latin
  Library (CLTK collection), used for the five-gram model of the cipher solvers; the Roman breviary (Divinum Officium),
  used in the tests on prayers and litanies.
\item \textbf{Handwritten gibberish:} the 38 released texts written by volunteers for \citet{gaskellbowern2022}.
\item \textbf{Published generators and ciphers:} the Naibbe cipher \citep{greshko2025}; the self-citation generator of
  \citet{timmschinner2020}; the generators U2 and U3 of \citet{whitehatnetizen2026}; the generator of
  \emph{voynich-fingerprint} and its procedure for a scribe working by hand \citep{sachak2026}.
\item \textbf{Medieval manuscripts transcribed at the level of letter forms}, with original lines and pages: six
  Nordic manuscripts from Menota \citep{menota} (AM 519 a 4to, Icelandic, c.~1280; AM 677 4to, Icelandic, c.~1200--1225;
  AM 60 4to, Norwegian, c.~1320; AM 242 fol, Icelandic, c.~1350; Holm A 10, Swedish, c.~1500--1550; AM 302 fol,
  Norwegian, c.~1300) and 73 German manuscripts of 1350--1500 from the Reference Corpus of Early New High German
  \citep{ref2021}. These are 79 manuscripts, not 79 known individual scribes: the files do not mark changes of hand (with one
exception in AM 677 4to).
\item \textbf{A historical cipher:} the Copiale, a mid-eighteenth-century German homophonic cipher of 105 pages, with
  the transcription and decipherment of \citet{knight2011,knight2012}.
\end{itemize}
Sources, versions and licences are listed in Appendix~\ref{app:sources}. All corpora except the Voynich
transliterations are kept outside the repository and must be downloaded from their sources.

% ---------------------------------------------------------------------------
\section{Methods}
\label{sec:methods}

\subsection{How the work was organised}

Every experiment followed the same cycle, each step in a separate commit of a version-controlled repository:
\begin{enumerate}
\item the program was tried only on synthetic texts and controls, never on the Voynich;
\item a preregistration was committed: question, data, measure and the criterion that would decide the outcome;
\item the program was committed;
\item the experiment was run by a script that records the code version, data fingerprints, random seeds and output;
\item results and a dated entry of the laboratory notebook, which is only ever appended to, were committed.
\end{enumerate}
The repository is public \citep{repository}. The preregistrations are internal: they are dated files committed before
each run, but they carry no third-party timestamp. Appendix~\ref{app:log} summarises the log of all 716 experiments
(e01--e3c92) with the outcome of each against its preregistered criterion; 698 were run, 670 were preregistered,
80 declare a deviation, and 19 of these are cases in which the author read an outcome differently from the literal criterion
(each listed with both readings).

\paragraph{Exploration and confirmation} The research programme was exploratory: each question arose from earlier
results on the same text. Confirmation on data not used to formulate a hypothesis is limited to a few cases: the
prediction that \eva{qo-} would be almost absent from drawing labels, made before looking at them (it was already known
from \citealt{zandbergen_labels}); the replications on the Takahashi and Glen Claston transliterations; and the
comparisons with the 73 German manuscripts, collected after the measures had been fixed on the Nordic ones.

\subsection{Controls}

\begin{itemize}
\item \textbf{Synthetic validation.} Every measure was first run on texts built with the effect present and absent. A
  measure that did not separate the two cases was not used; three experiments were dropped for this reason and declared
  as gaps (Section~\ref{sec:limits}).
\item \textbf{Positive and negative controls} in every decipherment test: a known text enciphered in the same way, which
  the method must solve, and a text in another language.
\item \textbf{Nulls that keep what the hypothesis does not concern}: for example, word order shuffled within the line,
  choices shuffled among occurrences of the same word, or (since version~2) nulls that keep word form or position in the
  line (Section~\ref{sec:rules}).
\end{itemize}

\subsection{Main measures}

The measures are described in words where they are used; their definitions are in Appendix~\ref{app:measures}.
\begin{itemize}
\item \textbf{Junction:} mutual information between the last glyph of a word and the first glyph of the next word in the
  same line, minus its mean under a null that shuffles the words of each line.
\item \textbf{Junction rule:} for a word that occurs in two forms (\eva{qo-}/\eva{o-}, \eva{-l}/\eva{-r}), the
  dependence of the form on the neighbouring word.
\item \textbf{Agreement of choices:} for binary choices between variant forms, a kappa-like agreement between nearby
  words, corrected for the bias of the expected value by subtracting the agreement obtained when the choices are
  shuffled among occurrences of the same word.
\item \textbf{Copying from the line above:} excess of identical or nearly identical words in the lines above over what
  the paragraph's vocabulary predicts.
\item \textbf{Left margin:} ratio of observed to expected identical line starts in consecutive lines, the expectation
  coming from shuffling the order of the lines within the paragraph.
\end{itemize}

\subsection{Uncertainty and multiplicity}

Intervals are bootstrap percentiles (2{,}000 resamples). In version~1 we resampled pages; but pages of the same
bifolium are not independent (Section~\ref{sec:book}), so the main intervals are now computed by resampling bifolia
(experiment \expcode{e3c86}). They are 0.84 to 1.56 times as wide as the page intervals (narrower in a few cases, as can
happen with about 50 groups), and each of the ten main
measures re-run in this way keeps its interval clear of the null value (copying from the line above and the junction
rules were not re-run); with only 16 quires, resampling quires widens most intervals further and one drift
(\eva{k}/\eva{t}) touches zero. For the main tests we declared seven families (one per finding) and applied Holm's
correction within each, and Bonferroni's correction over the 715 experiments logged before this test as a stricter reference
(\expcode{e3c92}): all 23 tests survive Holm, 15 survive Bonferroni (Table~\ref{tab:multiplicity} in
Appendix~\ref{app:measures}). Effect sizes are given with their intervals; $z$ values are given only as a guide.

\subsection{Use of generative AI}

The work was carried out in a continuous conversation with Claude, a large language model by Anthropic (mainly
Claude Opus 5.5), used through Claude Code in September and October 2026. Claude proposed experiments, wrote the
preregistrations, the analysis programs (Python 3.12) and the notebook entries, ran the experiments, summarised their
outputs and drafted the text of this paper. The voynichizer was developed with the same system in a separate
conversation on the same repository. The author set the questions and the direction of the research, chose among the
avenues proposed, approved every data download and every change of scope, and reviewed results and text. The research
conversation processed about 2.8 billion tokens, almost all of them the context of the conversation read again at each
step.
